% TOP_PROBLEM: {"id":30007524,"problem_number":"LOCAL-30007524","title":"Inflation-expectation formation","category_id":21,"category":{"id":21,"name":"economics","display_name":"Economics","description":"Open economic theory statements, published research agendas, and proposed scoped specifications; question provenance and historical priority are qualified per record.","slug":"economics","order_index":21,"created_at":"2026-10-08T00:00:00Z"},"status":"research_question","external_url":null,"legacy_ids":[],"published":false,"statement":"Identify and validate how personally experienced prices, information, and attention causally affect household inflation-belief updating and actual decisions.","economics":{"original_id":13,"field":"Inflation and monetary policy","kind":"identification","formalization_basis":"adapted_research_question","source_status":"explicit_open","priority_status":"earliest_found","priority_note":"This is the earliest explicit relevant statement verified in this audit, not a claim of original priority; older literature and earlier versions may contain antecedents.","earliest_verified_reference":{"authors":["Janet L. Yellen"],"title":"Macroeconomic Research After the Crisis","year":2016,"url":"https://www.federalreserve.gov/newsevents/speech/yellen20161014a.htm","doi":null,"locator":"Section “Inflation Dynamics,” paragraphs on whose expectations matter and how expectations form","evidence_summary":"The speech identifies unresolved questions about which agents’ inflation expectations matter, their quantitative influence, their formation, and the role of central-bank communication."},"current_reference":{"authors":["European Central Bank"],"title":"Research agenda: Forecasting, macroeconomic and financial analysis","year":null,"url":"https://www.ecb.europa.eu/pub/economic-research/research_agenda/forecasting/html/index.en.html","doi":null,"locator":"“Understanding the inflation process,” selected areas of focus","evidence_summary":"The agenda explicitly includes Phillips-curve determinants and their evolution, wage–price dynamics, and the measurement, determinants, and role of inflation expectations."},"formalization_note":"The source explicitly asks whose expectations matter and how they form. The transition model and experimental identification requirements are an adaptation."},"statusQualification":"Published question or research agenda verified at the cited locator. The mathematical specification is adapted for this catalogue and includes additional assumptions; source publication does not establish first-ever priority or certify this exact model remains unsolved. The source explicitly asks whose expectations matter and how they form. The transition model and experimental identification requirements are an adaptation.","statusReviewedAt":"2026-10-08"}
% FIRST_POSTED: 2026-10-08
% ENTRY_KIND: statement_only
% !TeX program = lualatex
\documentclass[11pt]{article}
\usepackage{fontspec}
\usepackage[a4paper,margin=25mm]{geometry}
\usepackage{amsmath,amssymb,mathtools}
\usepackage{xurl}
\usepackage[hidelinks]{hyperref}
\newcommand{\sref}[2]{\hyperlink{source:#1:#2}{[S#2]}}
\setlength{\emergencystretch}{3em}
\begin{document}
% problemId:30007524
\section*{Inflation-expectation formation}\label{30007524}
\subsection{Definitions and mathematical statement}
For household \(i\), let \(m_{it}\) denote its subjective mean of next-year inflation, \(p_{it}\) its personally experienced price change, \(n_{it}\) the inflation information it receives, and \(a_{it}\) measured attention. Specify the belief-transition model
\[
m_{i,t+1}=G_\theta(m_{it},p_{it},n_{it},a_{it},h_{it})+\epsilon_{i,t+1},
\]
where \(h_{it}\) contains predetermined household characteristics, \(\theta\) is a finite-dimensional parameter, and \(\epsilon_{i,t+1}\) is a mean-zero innovation conditional on the stated information. The experimental design independently randomizes informative messages and attention incentives; personally experienced prices require a separate valid instrument or an explicit partial-identification strategy.
Identify the causal effects of information content, exposure to personally salient prices, and attention on both the mean and dispersion of subjective inflation distributions. Test whether one estimated transition rule transports to different inflation regimes and communication formats, using held-out households and occasions. Distinguish persistent updating from temporary survey-response priming and reconcile forecasts at different horizons. To establish macroeconomic relevance, connect beliefs to actual consumption or pricing through independently identified outcome equations rather than assuming that survey revisions imply behavior. The target is an empirically disciplined formation mechanism and its policy responsiveness within a declared population; unrestricted household-specific updating rules do not constitute a solution.

\par\textbf{Formulation and scope.} The source explicitly asks whose expectations matter and how they form. The transition model and experimental identification requirements are an adaptation.

\par\textbf{Open-question provenance.} An explicit question or research agenda was verified at the source locator. This does not by itself certify that every later version remains unresolved \sref{30007524}{1}.

\par\textbf{Historical priority.} This is the earliest explicit relevant statement verified in this audit, not a claim of original priority; older literature and earlier versions may contain antecedents.

\subsection{Short English statement}
Identify and validate how personally experienced prices, information, and attention causally affect household inflation-belief updating and actual decisions.

\subsection{Sources}
\begin{itemize}\raggedright
\item[\textnormal{[S1]}]\hypertarget{source:30007524:1}{} Janet L. Yellen. 2016. Macroeconomic Research After the Crisis. Section “Inflation Dynamics,” paragraphs on whose expectations matter and how expectations form.
\url{https://www.federalreserve.gov/newsevents/speech/yellen20161014a.htm}.
{\small\textit{Earliest verified question/agenda reference; historical priority qualified below. The speech identifies unresolved questions about which agents’ inflation expectations matter, their quantitative influence, their formation, and the role of central-bank communication.}}

\item[\textnormal{[S2]}]\hypertarget{source:30007524:2}{} European Central Bank. Year not verified. Research agenda: Forecasting, macroeconomic and financial analysis. “Understanding the inflation process,” selected areas of focus.
\url{https://www.ecb.europa.eu/pub/economic-research/research_agenda/forecasting/html/index.en.html}.
{\small\textit{Supporting or current literature; see evidence qualification. The agenda explicitly includes Phillips-curve determinants and their evolution, wage–price dynamics, and the measurement, determinants, and role of inflation expectations.}}
\end{itemize}
\end{document}
